\documentclass{sciposter}
%Estes são os pacotes utilizados no arquivo
\usepackage{lipsum}
\usepackage{setspace}
\usepackage{epsfig}
\usepackage{amsmath}
\usepackage{amssymb}
\usepackage{multicol}
\usepackage{graphicx,url}
\usepackage[table,xcdraw]{xcolor}
\usepackage{bm}
\usepackage{color,colortbl,multirow}
\usepackage[brazil]{babel}  
\usepackage[utf8]{inputenc}
\usepackage[T1]{fontenc}
\usepackage{helvet}

%Esses são álguns dos teoremas, definições, proposições, que poderá usar ao longo do arquivo. Poderá criar outros se for necessário.
\newtheorem{Def}{Definition}
\newtheorem{teo1}{Teorema}
\newtheorem{teo}{Teorema}[section]
\newtheorem{id}[teo]{Ideia da Demonstra\c c\~ao}
\newtheorem{propr}[teo]{Propriedades}
\newtheorem{lem}[teo]{Lemma}
\newtheorem{coro}[teo]{Corollary}
\newtheorem{prop}[teo]{Proposi\c c\~ao}
\newtheorem{obs}[teo]{Observação}%Remark}
\newtheorem{exe}[teo]{Exemplo}
\newtheorem{defin}[teo]{Defini\c c\~ao}
\newtheorem{af}[teo]{Afirma\c c\~ao}
\newtheorem{ques}[teo]{Quest\~ao}
%________________________________________________________________
%Esquema de cores, para os tópicos.
\definecolor{BoxCol}{rgb}{0, 0, 0}
\definecolor{pretomostra}{rgb}{1, 1, 1}
\definecolor{SectionCol}{rgb}{1, 1, 1}
%_________________________________________________________________

%Título do projeto
\title{Título do Trabalho}

%Autor, Orientador e Coautores
\author{Nome do autor apresentador\footnote{Graduando da Instituição de Ensino Superior},\\ Nome do orientador\footnote{Orientador do projeto e professor da Instituição de Ensino Superior}}

%Instituição
\institute{Instituição de Ensino Superior}

\email{email@gmail.com}


\begin{document}

\begin{picture}(2cm,2cm)
  \put(1750,-330){\includegraphics[scale=0.55]{Logo Mostra IC.png}} %Logo universidade Federal de Ouro Preto
\end{picture}
\vspace{-3cm}
\begin{picture}(1.5cm, 1.5cm)
  \put(-40,-100){\includegraphics[scale=0.5]{logo ufop.png}} %Logo II Mostra de IC
\end{picture}

\begin{picture}(2cm,2cm)
  \put(5,-375){\includegraphics[scale=0.4]{Logo Demat.png}} %Logo DEMAT - UFOP
\end{picture}

\maketitle

\begin{multicols*}{2}

%________________________________________________________________
%\setlength{\columnseprule}{0pt} %Linha entre as colunas
%_________________________________________________________________

%%% Introdução
\section{\large{Seção 1}} %Você pode alterar o nome das seções
\vspace{0.2cm} 

\lipsum[10] %Texto em Latim (Pode ser apagado do arquivo)

%_______________________________________________________________
\section{\large{Seção 2}} %Você pode alterar o nome das seções
\vspace{0.2cm}
\lipsum[15] %Texto em Latim (Pode ser apagado do arquivo)

\begin{teo}
    Este é um exemplo de Teorema
    \label{teo1}
\end{teo}

\begin{figure}[H]
    \centering
    \includegraphics[width=10cm]{Logo Mostra IC.png}
    \caption{Este é um exemplo de figura}
    \label{fig2.1}
\end{figure}

Este é um exemplo de equação:
\begin{equation}
    F(x) = \int_a ^x f(t)dt
    \label{eq:1}
\end{equation}

\begin{figure} 
    \label{fig:2}
    \centering
    \subfigure{\includegraphics[scale=0.5]{logo ufop.png}} \qquad
    \subfigure{\includegraphics[scale=0.5]{Logo Demat.png}}
    \label{fig:2b}
\caption{Este é um outro exemplo de figura}
\end{figure}

\lipsum %Texto em Latim (Pode ser apagado do arquivo)

\begin{itemize}
    \item Este;
    \item É;
    \item Um;
    \item Exemplo;
    \item De;
    \item Lista.
\end{itemize}

\lipsum[20] %Texto em Latim (Pode ser apagado do arquivo)

\vspace{1.2cm}
Para referênciar algum teorema, imagem ou equação, usar \ref{teo1}, \ref{fig:2}, \ref{fig:2b} e \ref{eq:1}.

Para citações, usar \cite{burgos2016solution}, \cite{Eder}, \cite{Leibniz} e \cite{de2001edp}.
%__________________________________________________________________
\section{\large{Seção 3}} %Você pode alterar o nome das seções
\vspace{0.2cm}

\lipsum[5] %Texto em Latim

\section*{\large{Seção 1}} %Você pode alterar o nome das seções; Observe que o asterisco neste comando, evita que esta seção seja numerada.
\vspace{0.2cm}

Agradeço à agência que financiou este trabalho.

\nocite{*}   %Este comando garante que as referências aparecerão, mesmo que não sejam referenciadas no pôster.

%As referencias encontram-se no arquivo ref.bib

\bibliographystyle{plain} % Estilo de citação
\bibliography{ref}  %nome do arquivo bib, com as bibliografias

\end{multicols*}

\end{document}